\chapter{Bread, Water, Wine}
\label{ch:food}

Food at last --- and three words that sit at the centre of the French table.
One of them hides the loveliest word-story in English; another is the most
worn-down word French ever took from Latin. \emph{Practice lessons:
\texttt{lessons/FR-C11-*.md}.}

\section{\emph{le pain} --- bread, and the people you share it with}
\label{lesson:pain}

\textbf{le pain} (``bread'') is the heart of the French table --- and its Latin
root hides one of the loveliest word-stories in English: your
\textbf{companion} is, literally, someone you \textbf{share bread} with.

\begin{sounds}
\emph{le pain} carries the nasal vowel spelled \emph{-ain}: the air goes
through the nose, and the final \emph{-n} is not pronounced as a separate
consonant. It is masculine: \emph{le} pain.
\end{sounds}

\begin{cousinweb}
\textbf{le pain} (``bread'') $\leftarrow$ Latin \emph{p\=anis}. The root
\emph{p\=an-} feeds a surprising table of English words:\\[3pt]
\textbf{companion} / \textbf{company} $\leftarrow$ \emph{com-} (``with'') $+$
\emph{p\=anis} --- ``one you break bread with.'' A \emph{company} was first a
group sharing meals.\\[3pt]
\textbf{pantry} --- the room where the bread (and the rest of the food) is
kept.\\[3pt]
\textbf{pannier} --- a bread-basket, now the bag slung over a bicycle or a
donkey.
\end{cousinweb}

\begin{culture}
So \emph{le pain} is the same root as \textbf{companion}: French kept the
bread, English kept the friendship built on it. Two useful everyday forms:
\emph{un pain} (``a loaf'') and \emph{du pain} (``some bread'').
\end{culture}

\section{\emph{l'eau}, \emph{le vin} --- water and wine}
\label{lesson:eau-vin}

The two drinks of the French table. One of them, \textbf{eau}, is the single
most \textbf{worn-down} word French took from Latin: it began as \emph{aqua}
and ended as a bare ``oh.''

\begin{center}
\begin{tabular}{@{}ll@{}}
\toprule
French & English \\
\midrule
\emph{l'eau} & the water \\
\emph{le vin} & the wine \\
\bottomrule
\end{tabular}
\end{center}

\begin{sounds}
\emph{l'eau} is pronounced simply ``oh'' --- three letters, \emph{e-a-u},
spelling one vowel sound. It takes \emph{l'} rather than \emph{le} or \emph{la}
because it begins with a vowel. \emph{le vin} has the nasal vowel of
\emph{-in}.
\end{sounds}

\begin{cousinweb}
\textbf{l'eau} (``water'') $\leftarrow$ Latin \emph{aqua}. Watch the erosion:
\emph{aqua} $\to$ \emph{ewe} $\to$ \emph{eaue} $\to$ \emph{eau}, until only the
vowel ``oh'' remains. French quite literally \emph{drank the consonants out of}
\emph{aqua}.\\[3pt]
The original \emph{aqua} is still loud in English, because English borrowed the
Latin word \emph{whole}: \textbf{aqua}tic, \textbf{aqua}rium,
\textbf{aque}duct, \textbf{aqua}lung --- and Italian \emph{acqua} barely
changed at all. So \emph{l'eau} and \emph{aquarium} are the far ends of the
same word.\\[3pt]
\textbf{le vin} (``wine'') $\leftarrow$ Latin \emph{v\={\i}num}. Unlike
\emph{eau}, this one held its shape, and English shares it: \textbf{wine}
itself is \emph{v\={\i}num} borrowed into old Germanic, plus \textbf{vine},
\textbf{vin}egar (\emph{vin} $+$ \emph{aigre}, ``sour wine''),
\textbf{vint}age, \textbf{vine}yard, \textbf{vin}yl. Wine mattered enough to
keep its name across every border.
\end{cousinweb}
